\begin{lemma}[Residue independence for honest vector bundles]\label{rl-0011}\label{lem:independence-for-honest-vector-bundles}
  Let \(C\) be a \(T\)-stack over \(\mathcal Z\) admitting \(T\)-equivariant closed embeddings \(j_1:C\hookrightarrow E_1\) and \(j_2:C\hookrightarrow E_2\) into two different \(T\)-equivariant finite rank moving vector bundles over \(\mathcal Z\). Then for all \(\beta \in A^T_*(C)\),
  \[S_{(j_1,E_1)}(\beta) = S_{(j_2,E_2)}(\beta).\]
  We can therefore drop the reference to the embedding data and write \(S_C(\beta)\) without ambiguity.
\end{lemma}
\begin{proof}
  The strategy is to embed \(C\) into \(E_1\oplus E_2\) in two ways: first via
  \begin{align*}
    i_0:C\overset{j_1}{\hookrightarrow} E_1 = E_1\oplus 0 \overset{\iota}{\hookrightarrow} E_1 \oplus E_2,
  \end{align*}
  and then via the diagonal embedding
  \begin{align*}
    i_1:C\longhookrightarrow E_1\oplus E_2, \quad c\mapsto (j_1(c), j_2(c)).
  \end{align*}
  First define the interpolation map
  \begin{align*}
    \varphi:C\times \mathbb A^1 \to E_1\oplus E_2 \times \mathbb A^1; \quad \varphi(c, t) = (j_1(c), t\cdot j_2(c), t),
  \end{align*}
  noting that the fibers of \(\varphi\) at \(C\times\{0\}\) and \(C\times \{1\}\) recover \(i_0\) and \(i_1\) respectively. This map is \(T\)-equivariant if we let \(T\) act trivially on \(\mathbb A^1\).

  Now let \(\pr_C:C\times \mathbb A^1 \to C\) and \(\pr_{\mathcal Z}:\mathcal Z\times \mathbb A^1\to \mathbb Z\) be the projection maps and for \(\beta \in A^T_*(C)\) set
  \begin{align*}
    \Theta(\beta) :&= 0^!_{(E_1\oplus E_2) \times \mathbb A^1} (\varphi_*\pr^*_C\beta) \\
    &= 0^!_{(E_1\oplus E_2) \times \mathbb A^1} (\varphi_*(\beta \times [\mathbb A^1])) \in A^T_*(\mathcal Z\times\mathbb A^1).
  \end{align*}
  For each \(t\), the inclusion \(r_t:\{t\}\hookrightarrow \mathbb A^1\) yields Kresch's refined Gysin map \(r^!_t\). It acts a map \(r^!_t:A^T_*(Y)\to A^T_*(Y_t)\) for any \(Y\) over \(\mathbb A^1\), it commutes projective pushforward and flat pullback and \(r^!_t\circ \pr_C = \id\) on \(A^T_*(C)\). Therefore,
  \begin{align*}
    r^!_t\Theta(\beta) &= r^!_t0^!_{(E_1\oplus E_2) \times \mathbb A^1}(\varphi_*(\beta \times [\mathbb A^1])) \\
    &= 0^!_{E_1\oplus E_2}(r^!_t\varphi_{*}(\beta \times [\mathbb A^1])) \\
    &= 0^!_{E_1\oplus E_2}(\varphi_{t*}(\beta)).
  \end{align*}
  Since \(r^!_t\Theta(\beta) = (\pr_C)^{-1}\circ \Theta(\beta)\) is independent of \(t\), so is \(0^!_{E_1\oplus E_2}(\varphi_{t*}(\beta))\), and hence
  \begin{align*}
    0^!_{E_1\oplus E_2} (i_1\beta) = 0^!_{E_1\oplus E_2}(\varphi_{1*}(\beta)) = 0^!_{E_1\oplus E_2}(\varphi_{0*}(\beta)) = 0^!_{E_1\oplus E_2}(i_0\beta).
  \end{align*}

  All that is left is to transport the inverse Euler class through the various morphisms. Let \(f:E_1\oplus E_2\to \mathcal Z\) be the bundle map and factor it as \(f = g\circ h\) where \(h:E_1\oplus E_2\to E_1\) is the first projection and \(g:E_1\to \mathcal Z\) is the bundle map of \(E_1\). This means
  \begin{align*}
    0^!_{E_1\oplus E_2} = (f^*)^{-1} = (g^{*})^{-1}\circ (h^*)^{-1} = 0^!_{E_1} \circ (h^*)^{-1}.
  \end{align*}
  Now remember that we defined \(\iota:E_1\to E_1\oplus E_2\), so the self-intersection formula for a class \(\zeta\) in \(A_*^T(E_1)\) reads
  \begin{align*}
    (h^*)^{-1}(\iota_* \zeta) = e_T(g^*E_2) \cap \zeta,
  \end{align*}
  which we apply below to \(\zeta = j_{1*}\beta\).

  Applying the Whitney sum formula and the projection formula (\cite[\S 3.6]{kresch_CycleGroupsArtin1999}) then yields
  \begin{align*}
    S_{(i_1, E_1\oplus E_2)}(\beta) &= e^{-1}_T(E_1\oplus E_2) \cap 0^!_{E_1\oplus E_2}(i_{1*}\beta) \\
    &= e^{-1}_T(E_1)\cdot e_T^{-1}(E_2) \cap 0^!_{E_1\oplus E_2}(i_{0*}\beta) \\
    &= e^{-1}_T(E_1)\cdot e_T^{-1}(E_2) \cap 0^!_{E_1}\big((h^*)^{-1}(\iota_*j_{1*}\beta)\big) \\
    &= e^{-1}_T(E_1)\cdot e_T^{-1}(E_2) \cap 0^!_{E_1}\big(e_T(g^*E_2) \cap j_{1*}\beta\big) \\
    &= e^{-1}_T(E_1)\cdot e_T^{-1}(E_2)\cdot e_T(E_2) \cap 0^!_{E_1}\big(j_{1*}\beta\big) \\
    &= e_T^{-1}(E_1) \cap 0^!_{E_1}\big(j_{1*}\beta \big) \\
    &= S_{(j_1, E_1)}(\beta).
  \end{align*}
  Swapping the roles of \(E_1\) and \(E_2\) and using the interpolation
  \begin{align*}
    \varphi'(c,t) = (t·j1(c), j2(c), t)
  \end{align*}
  gives us \(S_{i_1, E_1\oplus E_2}(\beta) = S_{j_2, E_2}(\beta)\), and we are done.
\end{proof}
